\documentstyle[%


righttag%


,11pt%


,amssymb%


,russian%


,izvru%


]{amsart}





\newcommand{\tl}[3]{\medskip\noindent{\bf#1}\ #2 \dotfill #3 \par} 


\pagestyle{empty}


% \pagestyle{plain}


\begin{document}


\par


\vskip 2cm


\par


\bigskip


\bigskip


\bigskip


\par


\centerline{\Large СОДЕРЖАНИЕ}


\bigskip


\bigskip


%%%%%%%%%% Use \newline %%%%%%%%%%%








\tl{Алексеев Д.В.} 


{Приближение функций нескольких переменных


с весом\newline Чебышева--Эрмита}{3}


\tl{Коняев Ю.А.} 


{Метод расщепления в 


теории регулярных и сингулярных возмущений}{10}


\tl{Путилина А.В.}


{Асимптотика функции объема множества меньших


значений\newline полинома в $\Bbb R^n$}{16}


\tl{Сатаров Ж.С.}  


{Определяющие соотношения обобщенных ортогональных групп\newline над 


коммутативными локальными кольцами без единицы}{24}


\tl{Сумин М.И.} 


{Субоптимальное управление полулинейными  


эллиптическими уравнениями\newline с фазовыми  


ограничениями, I: принцип максимума для 


минимизирующих\newline последовательностей, 


нормальность}{33}


\tl{Титов С.С.}


{Решение нелинейных уравнений 


в аналитических полиалгебрах. II\ }{45}


\tl{Тронин С.Н., Копп О.А.}


{Матричные линейные операды}{53}


\tl{Шпирко С.В.}


{Решение игр многих лиц с помощью градиентного метода прогнозного \newline


типа}{63}


\bigskip


\bigskip





\centerline{КРАТКИЕ СООБЩЕНИЯ}





\bigskip


\bigskip





\tl{Денисюк И.Т.} 


{Одна задача сопряжения аналитических функций в аффинно \newline


преобразованных областях с кусочно-гладкими границами}{70}


\tl{Кадиев Р.И.}


{Достаточные условия устойчивости по части переменных


линейных\newline стохастических систем с последействием}{75}


\tl{Ковалев Л.В.}


{О внутренних радиусах


симметричных неналегающих областей}{80}








\vfill


\noindent\raisebox{2pt}{\copyright} Казанский госудаpственный унивеpситет


\end{document}


